\documentclass[11pt]{article}

% ======================================================================
% Comprehensive report for the SAVED-MODEL results:
%   bruise_colab_saved_eval.ipynb  +  bruise_colab_saved_analysis.ipynb
% Numbers from the two result zips; 22 figures from analysis/LATEST_IMAGES/BEST_RESULTS.
% Compile:  tools/tectonic/tectonic.exe docs/saved_models_report.tex
% Self-contained (standard article class; only CTAN packages).
% ======================================================================
\usepackage[margin=1in]{geometry}
\usepackage{booktabs}
\usepackage{amsmath,amssymb}
\usepackage{graphicx}
\usepackage{xcolor}
\usepackage{caption}
\usepackage{float}
\usepackage{enumitem}
\usepackage[hidelinks]{hyperref}
\usepackage{titlesec}

\graphicspath{{figures_saved/}}
\captionsetup{font=small,labelfont=bf}
\newcommand{\best}[1]{\textbf{#1}}
\newcommand{\sig}{\textcolor{red!70!black}{\textbf{sig.}}}
\newcommand{\ns}{n.s.}
\titleformat{\section}{\normalfont\Large\bfseries}{\thesection}{0.6em}{}
\setlength{\parskip}{0.4em}\setlength{\parindent}{0pt}

\title{\bfseries Saved-Model Evaluation and Analysis:\\
Five Bruise-Segmentation Models on the 28-Subject Consensus Test Set}
\author{Injury Analytics Lab --- Bruise Segmentation Project}
\date{Generated from \texttt{bruise\_colab\_saved\_eval.ipynb} and
\texttt{bruise\_colab\_saved\_analysis.ipynb}\\
(the project's original saved checkpoints; results in
\texttt{results\_final/*saved*.zip})}

\begin{document}
\maketitle

\begin{abstract}
\noindent
This report documents the evaluation and full visual analysis of the project's five
\emph{saved} segmentation models---a SegFormer-B2 teacher, SegFormer-B0 direct and
distilled students, and two YOLO26n-sem models---loaded as trained checkpoints (no
retraining) and scored on the fixed \textbf{185-image / 28-subject} test set against the
three-annotator \textbf{2-of-3 majority} target. Each model's decision threshold is fitted
on the 134-image validation set with a 1-D logit sweep (median of the 1-SE tie band). The
SegFormer models cluster at $0.773$--$0.786$ mean Dice with a \textbf{0--0.5\% complete-miss
rate}; here the \textbf{distilled B0 (0.786) is the top scorer}, edging the B2 teacher
(0.773). The YOLO models reach a comparable median ($\sim$0.85 native) but a much lower mean
(0.69--0.74) and a \textbf{7--10\% miss rate}. As with every analysis in this project: (i)
model-vs-model Dice differences are not resolvable at $n=28$ (distillation $+0.009$,
student$-$teacher $+0.013$, SegFormer$-$YOLO $+0.051$ all \ns); (ii) the resolvable axis is
the complete-miss rate; (iii) the three human experts agree with one another at only
\textbf{0.639} mean Dice and every model beats that ceiling---indeed all three SegFormer
models match the consensus \emph{significantly better than the annotator (Paul) they were
trained on} ($+0.07$--$0.09$, $P{\ge}0.99$); and (iv) an apparent light-skin deficit is
partly a lesion-size artefact but a residual light-skin \emph{recall} deficit survives size
and subject clustering for most models. Twenty-two figures are included.
\end{abstract}

\tableofcontents
\bigskip

% ======================================================================
\section{Scope and Setup}
% ======================================================================
These are the five \emph{already-trained} models saved in the project's local directories
(the checkpoints used across prior presentations), loaded and scored \textbf{without any
retraining}. The evaluation protocol is identical to the rest of the project:
\begin{itemize}[leftmargin=1.4em]
\item \textbf{Test set:} 185 white-light images from 28 subjects; target = per-pixel 2-of-3
      majority vote of three experts (Paul, Gbarimah, Erik). Scored on a $640\times640$ grid.
\item \textbf{Threshold:} fitted on the 134-image \emph{validation} set only, by sweeping a
      1-D cut on the raw bruise logit and choosing the median of all cuts within one standard
      error of the peak (a stable ``tie-band'' rule; Figure~\ref{fig:sweep}).
\item \textbf{YOLO} is scored two ways: \emph{native argmax} (its home-turf letterbox path)
      and \emph{custom /255} (the same 640 geometry SegFormer sees). Unless noted, the
      analysis uses native argmax as the core YOLO number.
\item \textbf{Statistics:} the 185 images are only 28 subjects, so all intervals are
      \textbf{subject-level cluster bootstraps} (resample subjects, $B=4000$); model-vs-model
      contrasts are \emph{paired}. \emph{These are single saved checkpoints, so there is no
      training-seed spread---the intervals reflect test-set (subject) sampling only.}
\end{itemize}

\begin{figure}[H]
\centering\includegraphics[width=\linewidth]{f1.png}
\caption{Validation threshold sweep for each model (mean Dice vs.\ the logit cut). The flat
top is the point: over a wide band of cuts the threshold is underdetermined, so we take the
median of the 1-SE tie band (dashed line) rather than the noisy argmax (dot).}
\label{fig:sweep}
\end{figure}

% ======================================================================
\section{Accuracy and Distribution}
% ======================================================================
Table~\ref{tab:acc} is the master result table (from \texttt{saved\_eval}); the YOLO rows
there are the custom-/255 path. Figure~\ref{fig:headline} gives the headline Dice
(YOLO shown at native argmax, its stronger median).

\begin{table}[H]
\centering
\caption{Master results: threshold fitted on val, scored on test (185 images). YOLO rows are
the custom-/255 path. Latency is single-frame 640$\to$640 on the GPU.}
\label{tab:acc}
\begin{tabular}{lccccccc}
\toprule
Model & Params & Thr. & Median Dice & Mean Dice & Miss \% & Median ms & FPS \\
\midrule
SegFormer-B2 teacher      & 27.35M & 0.674 & 0.834 & 0.773 & \best{0.00} & 32.3 & 30.8 \\
SegFormer-B0 direct       & 3.71M  & 0.574 & 0.813 & 0.777 & 0.54 & 15.9 & 62.6 \\
SegFormer-B0 distilled    & 3.71M  & 0.562 & 0.828 & \best{0.786} & 0.54 & 15.9 & 62.6 \\
YOLO26n direct (custom)   & 1.63M  & 0.182 & 0.825 & 0.699 & 9.19 & 7.9 & 125.2 \\
YOLO26n distilled (custom)& 1.63M  & 0.214 & 0.805 & 0.692 & 5.95 & 7.9 & 126.8 \\
\bottomrule
\end{tabular}
\end{table}

\begin{figure}[H]
\centering\includegraphics[width=0.82\linewidth]{f5.png}
\caption{Headline Dice (YOLO at native argmax). The \textbf{distilled B0 is the top mean
scorer (0.786)}, ahead of the B2 teacher (0.773)---i.e.\ on these saved checkpoints the
``student beats teacher'' pattern holds (it did \emph{not} survive the batch-equalised final
retraining, but it does here).}
\label{fig:headline}
\end{figure}

\begin{figure}[H]
\centering\includegraphics[width=\linewidth]{f6.png}
\caption{Left: per-image Dice distribution (white = median, red = mean); the long low-Dice
tail is why we report the median and treat failures separately. Right: survival curves
$P(\text{Dice}\ge x)$---the YOLO curves drop earliest at low $x$, i.e.\ more near-zero images.}
\label{fig:violin}
\end{figure}

\begin{figure}[H]
\centering\includegraphics[width=\linewidth]{f7.png}
\caption{Left: per-image precision vs.\ recall (big dot = mean). Right: IoU vs.\ Dice---a
monotone consistency check. YOLO tends to sit high-precision / lower-recall (clean but
under-firing), which foreshadows its miss and fairness behaviour.}
\label{fig:pr}
\end{figure}

% ======================================================================
\section{Safety and Failure Modes}
% ======================================================================
The decisive failure for an injury-documentation tool is a \textbf{complete miss}: a real
bruise, empty mask. Figure~\ref{fig:miss} shows the miss rates and the ``best Dice $\ne$
safest model'' effect---the YOLO models have the highest \emph{median} Dice yet by far the
highest miss rate.

\begin{figure}[H]
\centering\includegraphics[width=\linewidth]{f8.png}
\caption{Complete-miss rate with subject-bootstrap intervals (left) and miss vs.\ median
Dice (right). SegFormer-B2 never blanks (0\%); the YOLO models blank 7--10\% of images
despite the best medians.}
\label{fig:miss}
\end{figure}

\begin{figure}[H]
\centering\includegraphics[width=0.9\linewidth]{f9.png}
\caption{Over- vs.\ under-segmentation: predicted$\div$true bruise area (black dash =
median). All models sit slightly below 1 (mild under-segmentation); the YOLO models carry a
spike of images at ratio 0 (the blank masks).}
\label{fig:overunder}
\end{figure}

\begin{figure}[H]
\centering\includegraphics[width=0.9\linewidth]{f10.png}
\caption{YOLO's two evaluation paths are not interchangeable: native argmax vs.\ custom
/255 per image. Same weights, different eval---report both and state which.}
\label{fig:paths}
\end{figure}

% ======================================================================
\section{Which Differences Are Real?}
% ======================================================================
Every marginal interval (Figure~\ref{fig:ci}) overlaps heavily; the paired contrasts
(Table~\ref{tab:paired}, Figure~\ref{fig:forest}) confirm that \textbf{only the miss-rate
gap clears zero}. Notably, on these saved checkpoints even SegFormer$-$YOLO on Dice is
\emph{not} significant ($+0.051$, CI spans zero), unlike the batch-equalised final run.

\begin{table}[H]
\centering
\caption{Paired subject-level cluster-bootstrap contrasts ($B=4000$). Dice rows except the
last (complete-miss rate, percentage points).}
\label{tab:paired}
\begin{tabular}{lcccc}
\toprule
Contrast & $\Delta$ & 95\% CI & $P(\Delta{>}0)$ & Verdict \\
\midrule
Distillation (B0-dist $-$ B0-direct) & $+0.009$ & $[-0.009,+0.029]$ & 0.82 & \ns \\
Student $-$ teacher (B0-dist $-$ B2) & $+0.013$ & $[-0.023,+0.052]$ & 0.74 & \ns \\
SegFormer $-$ YOLO (B0-dist $-$ YOLO-direct) & $+0.051$ & $[-0.011,+0.117]$ & 0.94 & \ns \\
YOLO distilled $-$ direct & $+0.006$ & $[-0.019,+0.030]$ & 0.67 & \ns \\
\textbf{MISS\%: YOLO-direct $-$ B0-dist} & $+9.19$ & $[+1.83,+17.54]$ & 1.00 & \sig \\
\bottomrule
\end{tabular}
\end{table}

\begin{figure}[H]
\centering\includegraphics[width=0.85\linewidth]{f11.png}
\caption{Marginal mean-Dice intervals. At 28 subjects the five models are hard to separate.}
\label{fig:ci}
\end{figure}

\begin{figure}[H]
\centering\includegraphics[width=0.9\linewidth]{f12.png}
\caption{Paired effect sizes; red = interval clears zero. Only the miss-rate contrast is
resolvable. $P(\Delta{>}0)$ distinguishes an underpowered null (distillation 82\%) from a
coin flip (student$-$teacher 74\%).}
\label{fig:forest}
\end{figure}

\begin{figure}[H]
\centering\includegraphics[width=\linewidth]{f13.png}
\caption{Left: per-subject mean Dice (hardest subjects on top); rows dark across every column
are hard \emph{subjects}, not model failures---and with only 28 rows, one or two of them move
the whole average. Right: Spearman correlation of per-image Dice; the models largely agree on
which images are hard.}
\label{fig:subject}
\end{figure}

% ======================================================================
\section{Fairness Across Skin Tone (ITA)}
% ======================================================================
\label{sec:fair}
Groups are ITA (objective, pixel-computed). This is \textbf{exploratory} at $n=28$
($\approx$9--17 subjects/group). SegFormer-B2 and B0-direct show a nominally significant
omnibus effect (Kruskal $p=0.016$, $0.027$); the striking pattern is YOLO's
\textbf{light-skin under-detection}: its complete-miss rate on light skin reaches
\textbf{28\%} (native direct) against 0--7\% elsewhere (Figure~\ref{fig:fairheat}, right).

\begin{figure}[H]
\centering\includegraphics[width=\linewidth]{f14.png}
\caption{Per-ITA-group median Dice (left), mean recall (middle), and complete-miss rate
(right) for every variant. The miss panel is the story: the YOLO variants concentrate their
blank masks on the lightest-skin group.}
\label{fig:fairheat}
\end{figure}

\begin{figure}[H]
\centering\includegraphics[width=\linewidth]{f15.png}
\caption{Left: best$-$worst-group median-Dice gap per variant (red = Kruskal $p<0.05$).
Right: YOLO-direct per-group precision vs.\ recall---light skin sits at high precision / very
low recall, i.e.\ \emph{under-detection}, not confusion.}
\label{fig:fairgap}
\end{figure}

\subsection{Is the light-skin gap just lesion size?}
Light-skin bruises are the \emph{smallest} in this cohort (Figure~\ref{fig:sizeita}), and
small bruises are harder for every model---so size is a confound. We test whether the signal
survives it with subject-clustered regressions (Table~\ref{tab:sizereg}) and stratification
(Figure~\ref{fig:e4}). \textbf{It largely survives:} the light coefficient on recall stays
significant with $\log(\text{size})$ controlled for the B2 teacher ($-0.195$, $p=0.003$),
B0-direct ($-0.098$, $p=0.016$), and both YOLO models ($-0.255$/$-0.246$, $p=0.014$/$0.003$);
B0-distilled is borderline ($p=0.057$). Within the smallest size tercile, YOLO-direct still
blanks \textbf{35\%} of light-skin images vs.\ 9.5\% of the rest. This is a stronger residual
signal than the final batch-equalised run showed---still a hypothesis at 9 light-skin
subjects, but a consistent one.

\begin{table}[H]
\centering
\caption{Light-skin coefficient with $\log_{10}(\text{size})$ controlled, subject-clustered
SEs. Negative = worse on light skin at equal size.}
\label{tab:sizereg}
\begin{tabular}{lccc}
\toprule
Model & Outcome & light coef & $p$ (clustered) \\
\midrule
SegFormer-B2 teacher   & recall & $-0.195$ & \best{0.003} \\
SegFormer-B0 direct    & recall & $-0.098$ & \best{0.016} \\
SegFormer-B0 distilled & recall & $-0.097$ & 0.057 \\
YOLO26n direct         & recall & $-0.255$ & \best{0.014} \\
YOLO26n distilled      & recall & $-0.246$ & \best{0.003} \\
YOLO26n direct         & miss (logit) & $+1.56$ & 0.088 \\
\bottomrule
\end{tabular}
\end{table}

\begin{figure}[H]
\centering\includegraphics[width=0.8\linewidth]{f18.png}
\caption{Bruise size by ITA group. Light skin has the smallest lesions---the confound that
Table~\ref{tab:sizereg} and Figure~\ref{fig:e4} control for.}
\label{fig:sizeita}
\end{figure}

\begin{figure}[H]
\centering\includegraphics[width=0.7\linewidth]{f19.png}
\caption{YOLO-direct complete-miss rate, light vs.\ rest, within each size tercile. The gap
does not vanish under size control---it is largest exactly where lesions are smallest.}
\label{fig:e4}
\end{figure}

% ======================================================================
\section{Bruise Size}
% ======================================================================
Small bruises are intrinsically harder (a few pixels of error cost proportionally more
Dice), and misses concentrate on them (Figure~\ref{fig:size}). Recall degrades on the
smallest lesions for the YOLO models while SegFormer stays flatter
(Figure~\ref{fig:sizeq}).

\begin{figure}[H]
\centering\includegraphics[width=\linewidth]{f16.png}
\caption{Bruise-size distribution (median 11{,}324\,px), Dice vs.\ size, and complete-miss
rate by size quartile. Misses concentrate on the smallest bruises for the YOLO models.}
\label{fig:size}
\end{figure}

\begin{figure}[H]
\centering\includegraphics[width=\linewidth]{f17.png}
\caption{Recall and precision by bruise-size quintile. Precision is roughly flat; recall is
where the smallest-lesion penalty shows up.}
\label{fig:sizeq}
\end{figure}

% ======================================================================
\section{The Annotation Ceiling}
% ======================================================================
The most important result concerns the labels. The three experts agree with one another at
only \textbf{0.639} mean Dice (Figure~\ref{fig:ceiling}); every model scores above that
against the consensus. Paul---whose masks the models train on---is the outlier, matching the
majority least (0.700) of the three.

\begin{figure}[H]
\centering\includegraphics[width=0.85\linewidth]{f20.png}
\caption{Inter-annotator agreement vs.\ model-vs-consensus. Every model beats the
human--human average (dashed, 0.639).}
\label{fig:ceiling}
\end{figure}

\subsection{The model beats the annotator it learned from}
Each model trains only on Paul's masks yet is scored against the three-expert majority.
Table~\ref{tab:paul} and Figure~\ref{fig:paul}: all three SegFormer models agree with the
consensus \emph{significantly better than Paul does}. This effect ($\approx+0.08$) is an
order of magnitude larger than any model-vs-model gap, and it is the study's only significant
accuracy result---the models are \textbf{label-limited, not capacity-limited}.

\begin{table}[H]
\centering
\caption{Model vs.\ Paul against the same majority target (paired subject bootstrap, $B=4000$).}
\label{tab:paul}
\begin{tabular}{lcccc}
\toprule
Model & Model vs.\ maj. & $\Delta$ (model $-$ Paul) & 95\% CI & $P(\Delta{>}0)$ \\
\midrule
SegFormer-B2 teacher   & 0.773 & $+0.073$ & $[+0.019,+0.131]$ & 0.998 \\
SegFormer-B0 direct    & 0.777 & $+0.077$ & $[+0.010,+0.140]$ & 0.988 \\
SegFormer-B0 distilled & 0.786 & $+0.087$ & $[+0.018,+0.156]$ & 0.994 \\
YOLO26n direct (native)& 0.735 & $+0.036$ & $[-0.028,+0.107]$ & 0.85 \\
YOLO26n distilled (native)& 0.741 & $+0.042$ & $[-0.022,+0.107]$ & 0.89 \\
\bottomrule
\end{tabular}
\end{table}

\begin{figure}[H]
\centering\includegraphics[width=0.85\linewidth]{f21.png}
\caption{$\Delta$ Dice (model $-$ Paul) vs.\ the same majority, with subject-bootstrap
intervals. All three SegFormer bars clear zero.}
\label{fig:paul}
\end{figure}

% ======================================================================
\section{Efficiency}
% ======================================================================
\begin{figure}[H]
\centering\includegraphics[width=0.8\linewidth]{f22.png}
\caption{Accuracy vs.\ speed (median Dice $\times$ benchmarked FPS). YOLO is far faster
($\sim$125 FPS, 1.6\,M params) but pays in miss rate; SegFormer-B0 is the balance point;
the B2 teacher is the accuracy anchor at $\sim$31 FPS.}
\label{fig:speed}
\end{figure}

% ======================================================================
\section{Qualitative Galleries}
% ======================================================================
\begin{figure}[H]
\centering\includegraphics[width=\linewidth]{f2.png}
\caption{Predictions across easy $\to$ hard images (green = truth, red = model). The hardest
row is a tiny, faint bruise that every model but the B2 teacher misses entirely.}
\label{fig:geasy}
\end{figure}

\begin{figure}[H]
\centering\includegraphics[width=\linewidth]{f3.png}
\caption{One representative image per ITA skin-tone group. Note the light-skin (top) row,
where YOLO-direct returns nothing (Dice 0.00) while SegFormer still localises the bruise.}
\label{fig:gita}
\end{figure}

\begin{figure}[H]
\centering\includegraphics[width=\linewidth]{f4.png}
\caption{One representative image per bruise-size quartile. The smallest (top) row is where
the YOLO models blank; all models handle the largest lesions well.}
\label{fig:gsize}
\end{figure}

% ======================================================================
\section{Discussion: how these saved models compare to the retrained run}
% ======================================================================
These are the project's \emph{original} checkpoints (one per model, no seed replication);
the batch-equalised, natively-retrained five-model run is reported separately in
\texttt{final\_run\_report.pdf}. The two agree on the structural findings and differ on two
points worth stating plainly:
\begin{enumerate}[leftmargin=1.4em]
\item \textbf{Student vs.\ teacher.} Here the distilled B0 (0.786) tops the B2 teacher
      (0.773); in the retrained run, once batch/step budget was equalised, the B2 led. Both
      gaps are within the annotation noise ($\Delta$ n.s.), so neither ordering is a real
      capability claim---but it is a clean example of how a confound (training budget) can
      flip a headline that the statistics never supported.
\item \textbf{Fairness signal.} On these saved models the residual light-skin recall deficit
      is significant (after size + clustering) for \emph{most} models including SegFormer,
      slightly stronger than the retrained run. At 9 light-skin subjects it remains a
      hypothesis, but a consistent one across both model sets.
\end{enumerate}
Everything else matches: SegFormer is essentially miss-free while YOLO is not; the miss rate,
not Dice, is the axis that separates architectures by more than label noise; and the
annotation ceiling (models beat the annotator they learned from) is the dominant, and only
significant, accuracy result.

% ======================================================================
\section{Reproducibility}
% ======================================================================
\begin{itemize}[leftmargin=1.4em]
\item \textbf{Evaluation:} \texttt{bruise\_colab\_saved\_eval.ipynb}
      (generator \texttt{scripts/46\_generate\_saved\_eval\_notebook.py}) --- loads the saved
      checkpoints via the \texttt{pipeline} package, sweeps the threshold on val, scores test,
      YOLO two ways, ITA fairness, annotation ceiling. Outputs: the \texttt{saved\_eval\_*} zip.
\item \textbf{Analysis / figures:} \texttt{bruise\_colab\_saved\_analysis.ipynb}
      (generator \texttt{scripts/47\_generate\_saved\_analysis\_notebook.py}) --- the 22
      figures reproduced here, plus the paired contrasts, size-conditioned regression, and
      galleries. Outputs: the \texttt{saved\_analysis\_*} zip.
\item All tabular numbers are read directly from those two zips; the figures are the
      notebooks' own exports (\texttt{analysis/LATEST\_IMAGES/BEST\_RESULTS/}). Nothing is
      hand-transcribed.
\end{itemize}

\end{document}
